% Options for packages loaded elsewhere
% Options for packages loaded elsewhere
\PassOptionsToPackage{unicode}{hyperref}
\PassOptionsToPackage{hyphens}{url}
\PassOptionsToPackage{dvipsnames,svgnames,x11names}{xcolor}
%
\documentclass[
  letterpaper,
  DIV=11,
  numbers=noendperiod]{scrartcl}
\usepackage{xcolor}
\usepackage[tmargin=30mm,bmargin=30mm]{geometry}
\usepackage{amsmath,amssymb}
\setcounter{secnumdepth}{-\maxdimen} % remove section numbering
\usepackage{iftex}
\ifPDFTeX
  \usepackage[T1]{fontenc}
  \usepackage[utf8]{inputenc}
  \usepackage{textcomp} % provide euro and other symbols
\else % if luatex or xetex
  \usepackage{unicode-math} % this also loads fontspec
  \defaultfontfeatures{Scale=MatchLowercase}
  \defaultfontfeatures[\rmfamily]{Ligatures=TeX,Scale=1}
\fi
\usepackage{lmodern}
\ifPDFTeX\else
  % xetex/luatex font selection
\fi
% Use upquote if available, for straight quotes in verbatim environments
\IfFileExists{upquote.sty}{\usepackage{upquote}}{}
\IfFileExists{microtype.sty}{% use microtype if available
  \usepackage[]{microtype}
  \UseMicrotypeSet[protrusion]{basicmath} % disable protrusion for tt fonts
}{}
\makeatletter
\@ifundefined{KOMAClassName}{% if non-KOMA class
  \IfFileExists{parskip.sty}{%
    \usepackage{parskip}
  }{% else
    \setlength{\parindent}{0pt}
    \setlength{\parskip}{6pt plus 2pt minus 1pt}}
}{% if KOMA class
  \KOMAoptions{parskip=half}}
\makeatother
% Make \paragraph and \subparagraph free-standing
\makeatletter
\ifx\paragraph\undefined\else
  \let\oldparagraph\paragraph
  \renewcommand{\paragraph}{
    \@ifstar
      \xxxParagraphStar
      \xxxParagraphNoStar
  }
  \newcommand{\xxxParagraphStar}[1]{\oldparagraph*{#1}\mbox{}}
  \newcommand{\xxxParagraphNoStar}[1]{\oldparagraph{#1}\mbox{}}
\fi
\ifx\subparagraph\undefined\else
  \let\oldsubparagraph\subparagraph
  \renewcommand{\subparagraph}{
    \@ifstar
      \xxxSubParagraphStar
      \xxxSubParagraphNoStar
  }
  \newcommand{\xxxSubParagraphStar}[1]{\oldsubparagraph*{#1}\mbox{}}
  \newcommand{\xxxSubParagraphNoStar}[1]{\oldsubparagraph{#1}\mbox{}}
\fi
\makeatother


\usepackage{longtable,booktabs,array}
\usepackage{calc} % for calculating minipage widths
% Correct order of tables after \paragraph or \subparagraph
\usepackage{etoolbox}
\makeatletter
\patchcmd\longtable{\par}{\if@noskipsec\mbox{}\fi\par}{}{}
\makeatother
% Allow footnotes in longtable head/foot
\IfFileExists{footnotehyper.sty}{\usepackage{footnotehyper}}{\usepackage{footnote}}
\makesavenoteenv{longtable}
\usepackage{graphicx}
\makeatletter
\newsavebox\pandoc@box
\newcommand*\pandocbounded[1]{% scales image to fit in text height/width
  \sbox\pandoc@box{#1}%
  \Gscale@div\@tempa{\textheight}{\dimexpr\ht\pandoc@box+\dp\pandoc@box\relax}%
  \Gscale@div\@tempb{\linewidth}{\wd\pandoc@box}%
  \ifdim\@tempb\p@<\@tempa\p@\let\@tempa\@tempb\fi% select the smaller of both
  \ifdim\@tempa\p@<\p@\scalebox{\@tempa}{\usebox\pandoc@box}%
  \else\usebox{\pandoc@box}%
  \fi%
}
% Set default figure placement to htbp
\def\fps@figure{htbp}
\makeatother





\setlength{\emergencystretch}{3em} % prevent overfull lines

\providecommand{\tightlist}{%
  \setlength{\itemsep}{0pt}\setlength{\parskip}{0pt}}



 


\usepackage{fvextra}
\DefineVerbatimEnvironment{Highlighting}{Verbatim}{breaklines,commandchars=\\\{\}}
\KOMAoption{captions}{tableheading}
\makeatletter
\@ifpackageloaded{tcolorbox}{}{\usepackage[skins,breakable]{tcolorbox}}
\@ifpackageloaded{fontawesome5}{}{\usepackage{fontawesome5}}
\definecolor{quarto-callout-color}{HTML}{909090}
\definecolor{quarto-callout-note-color}{HTML}{0758E5}
\definecolor{quarto-callout-important-color}{HTML}{CC1914}
\definecolor{quarto-callout-warning-color}{HTML}{EB9113}
\definecolor{quarto-callout-tip-color}{HTML}{00A047}
\definecolor{quarto-callout-caution-color}{HTML}{FC5300}
\definecolor{quarto-callout-color-frame}{HTML}{acacac}
\definecolor{quarto-callout-note-color-frame}{HTML}{4582ec}
\definecolor{quarto-callout-important-color-frame}{HTML}{d9534f}
\definecolor{quarto-callout-warning-color-frame}{HTML}{f0ad4e}
\definecolor{quarto-callout-tip-color-frame}{HTML}{02b875}
\definecolor{quarto-callout-caution-color-frame}{HTML}{fd7e14}
\makeatother
\makeatletter
\@ifpackageloaded{caption}{}{\usepackage{caption}}
\AtBeginDocument{%
\ifdefined\contentsname
  \renewcommand*\contentsname{Table of contents}
\else
  \newcommand\contentsname{Table of contents}
\fi
\ifdefined\listfigurename
  \renewcommand*\listfigurename{List of Figures}
\else
  \newcommand\listfigurename{List of Figures}
\fi
\ifdefined\listtablename
  \renewcommand*\listtablename{List of Tables}
\else
  \newcommand\listtablename{List of Tables}
\fi
\ifdefined\figurename
  \renewcommand*\figurename{Figure}
\else
  \newcommand\figurename{Figure}
\fi
\ifdefined\tablename
  \renewcommand*\tablename{Table}
\else
  \newcommand\tablename{Table}
\fi
}
\@ifpackageloaded{float}{}{\usepackage{float}}
\floatstyle{ruled}
\@ifundefined{c@chapter}{\newfloat{codelisting}{h}{lop}}{\newfloat{codelisting}{h}{lop}[chapter]}
\floatname{codelisting}{Listing}
\newcommand*\listoflistings{\listof{codelisting}{List of Listings}}
\makeatother
\makeatletter
\makeatother
\makeatletter
\@ifpackageloaded{caption}{}{\usepackage{caption}}
\@ifpackageloaded{subcaption}{}{\usepackage{subcaption}}
\makeatother
\usepackage{bookmark}
\IfFileExists{xurl.sty}{\usepackage{xurl}}{} % add URL line breaks if available
\urlstyle{same}
\hypersetup{
  pdftitle={Homework 4: Model Selecction},
  pdfauthor={Avinash Aravind, aa2476},
  colorlinks=true,
  linkcolor={blue},
  filecolor={Maroon},
  citecolor={Blue},
  urlcolor={Blue},
  pdfcreator={LaTeX via pandoc}}


\title{Homework 4: Model Selecction}
\usepackage{etoolbox}
\makeatletter
\providecommand{\subtitle}[1]{% add subtitle to \maketitle
  \apptocmd{\@title}{\par {\large #1 \par}}{}{}
}
\makeatother
\subtitle{BEE 4850/5850}
\author{Avinash Aravind, aa2476}
\date{2026-03-27}
\begin{document}
\maketitle

\RecustomVerbatimEnvironment{verbatim}{Verbatim}{
showspaces = false,
showtabs = false,
breaksymbolleft={},
breaklines
% Note: setting commandchars=\\\{\} here will cause an error
}

\renewcommand*\contentsname{Contents}
{
\hypersetup{linkcolor=}
\setcounter{tocdepth}{3}
\tableofcontents
}

\begin{tcolorbox}[enhanced jigsaw, left=2mm, breakable, leftrule=.75mm, colframe=quarto-callout-note-color-frame, toptitle=1mm, coltitle=black, toprule=.15mm, rightrule=.15mm, bottomtitle=1mm, titlerule=0mm, colbacktitle=quarto-callout-note-color!10!white, bottomrule=.15mm, title=\textcolor{quarto-callout-note-color}{\faInfo}\hspace{0.5em}{Note}, arc=.35mm, opacitybacktitle=0.6, opacityback=0, colback=white]

All of the code used in this homework is accessible in the hw04.ipynb
file in my GitHub repository.

\end{tcolorbox}

\newpage

\section{Problem 1: Cherry Blossoms and
Temperature}\label{problem-1-cherry-blossoms-and-temperature}

\subsubsection{Problem 1.1}\label{problem-1.1}

\begin{figure}[h]

\centering{

\includegraphics[width=4.5in,height=\textheight,keepaspectratio]{"figs/11.png"}

}

\caption{\label{fig-11}March Temperature Anomalies and Cherry Blossom
Bloom Date}

\end{figure}%

The linear fit suggests that increasing temprature anomalies lead to a
decrease in the date of year that the cherry blossom bloom peaks; that
is, the bloom occurs earlier in the year. The slope of fit is about
-7.8, which suggests that a 1 degree C increase in temperature anomaly
leads to a bloom more than one week earlier.

\subsubsection{Problem 1.2}\label{problem-1.2}

The mean-squared error for this fitted model on the in-sample data is
about 22.32. I was unfortunately having trouble implementing the
log-loss in my code; I will attempt to understand how to implement this
after the solutions are released.

\subsubsection{Problem 1.3}\label{problem-1.3}

Using 5-fold cross-validation, the mean-squred errors for each fold were
roughly similar to the in-sample data, with MSE values from my run of
34.68, 23.15, 20.46, 22.82, and 18.74. Some of the folds have a
lower/better MSE than the full dataset, and some have a higher/worse
MSE, but there are no major differences.

\newpage

\subsubsection{Problem 1.4}\label{problem-1.4}

This model is relatively generalizable. The 5-fold cross-validation did
not produce meaningfully worse MSE values than the in-sample model,
suggesting that the model was not learning much more information from
the full dataset that it could not learn from a much smaller subset.
Thus, this model should be similarly good at predicting unseen values as
it is at predicting in-sample values.

\newpage

\section{Problem 2: Back in Chicago}\label{problem-2-back-in-chicago}

\subsubsection{Problem 2.1}\label{problem-2.1}

\begin{figure}[h]

\centering{

\includegraphics[width=4.5in,height=\textheight,keepaspectratio]{"figs/21.png"}

}

\caption{\label{fig-21}Chicago Deaths vs.~Time with Linear Fit}

\end{figure}%

The linear regression coefficient for this fit is -1.76e-17, very close
to zero, and the intercept is 128.75. As visible in Figure~\ref{fig-21},
there is basically no change in the model predicted number of deaths
through time, predicting about 129 deaths throughout the time period.

\newpage

\subsubsection{Problem 2.2}\label{problem-2.2}

\begin{figure}[h]

\centering{

\includegraphics[width=4.5in,height=\textheight,keepaspectratio]{"figs/22.png"}

}

\caption{\label{fig-22}Chicago Deaths vs.~Temperature with Linear Fit}

\end{figure}%

The linear regression coefficient for this fit is -0.29, and the
intercept is 129.96. Unlike the time fit, there is a trend in
Figure~\ref{fig-22}, with the coefficient suggesting that the model
predicts a decrease of a bit less than 1 death per 3 degrees increase in
mean daily temperature. The scatterplot of data appears to follow this
decreasing trend; however, the model grossly underestimates the very
high deaths during the heat wave.

\newpage

\subsubsection{Problem 2.3}\label{problem-2.3}

\begin{figure}[h]

\centering{

\includegraphics[width=4.5in,height=\textheight,keepaspectratio]{"figs/23.png"}

}

\caption{\label{fig-23}Chicago Deaths vs.~Temperature with Linear Fit
incorporating PM25}

\end{figure}%

The linear regression coefficients for this fit are -0.37 and 0.215, and
the intercept is 129.29. When incorporating the median pm25
concentration, the model predicts a greater decrease in temperature for
a given decrease in . However, it also predicts an increase in deaths
with increasing pm25 concentration. The resulting slope in
Figure~\ref{fig-23} is thus very similar to the one in
Figure~\ref{fig-22}.

\newpage

\subsubsection{Problem 2.4}\label{problem-2.4}

For the model with temperature only, the 5-fold cross-validation MSE
scores are 164.43, 170.09, 178.88, 321.36, and 203.51. For the model
with both temperature and pm25 concentration, the MSE scores are 140.53,
170.13, 183.30, 150.67, and 122.85. The dataset has the previously
mentioned spike in deaths during the heat wave of 1995; random selection
should split up this data, but it may explain the anomalously high MSE
values for folds 4 and 5 of the temperature-only dataset if the heat
wave days ended up exclusively in that fold. The model with both
temperature and pm25 appears to perform better, with lower MSE scores,
although not by much.

\subsubsection{Problem 2.5}\label{problem-2.5}

I was unfortunately unable to calculate AIC.




\end{document}
